\documentclass[12pt]{article}
\usepackage{amsmath}
\usepackage{amssymb}
\usepackage{amsthm}
\usepackage{cancel}
\usepackage{enumitem}
\usepackage{esdiff}
\usepackage{graphicx}
\usepackage[version=4,textfontname=sffamily,mathfontname=mathsf]{mhchem}
\usepackage{multicol}
\usepackage{siunitx}
\usepackage{ulem}
% \usepackage{pgfplots}
\usepackage{wrapfig}
\usepackage{xcolor}

\newcommand{\e}[1]{e^{#1}}
\newcommand{\ei}[1]{e^{i\left(#1\right)}}
\newcommand{\E}[1]{\times 10^{#1}}
\newcommand{\tagref}[1]{\tag{\ref{#1}}}

\renewcommand\qedsymbol{TENA}

\title{
    Homework \#14
    \\  \small
    PHYS 4D: Modern Physics
    \\
    Problems from Krane's textbook
    }
\author{Donald Aingworth IV}
\date{April 20, 2026}

\begin{document}
\maketitle
\DeclareSIUnit{\amu}{amu}

\section{Questions, Chapter 11}
    \subsection{Question 2}
        How should Eq. 11.7 be modified to be valid for \ce{MgO} and \ce{BaO}?
        \begin{equation}
            \label{eq:11.7} \tag{11.7}
            E_{\rm coh} = \frac{1}{2}(B)(2N_A) = BN_A
        \end{equation}
        
        \subsubsection{Answer}
            Let's take a closer look at that value $B$. 
            It can be expanded.
            \begin{equation}
                B = -U(R_0) = \frac{\alpha e^2}{4\pi\varepsilon_0 R_0} \left( 1 - \frac{1}{n} \right)
            \end{equation}

            There are two values here to be explored.
            The first is the Madelung constant $\alpha$, which is dependant on the binding structure. 
            \ce{MgO} and \ce{BaO} are both fcc structures, so they have Madelung constant $\alpha = 1.7476$.
            The second is $n$, which comes from an exponent of the equilibrium separation determined from compressability data.
            For \ce{MgO}, $n = 7$.
            For \ce{BaO}, $n = 9.5$.

    \subsection{Question 3}
        From Figure 11.11, estimate the Einstein temperature for lead. 
        (Hint: Consider Eq. 11.11 when $T = T_E$.)

        \subsubsection{Answer}

    \subsection{Question 4}
        A graph of $C/T$ vs. $T^2$ for solid argon (similar to Figure 11.13) goes through the origin; that is, its y-intercept is zero. 
        Explain.

        \subsubsection{Answer}

    \subsection{Question 7}
        (a) Why does the electrical conductivity of a metal decrease as the temperature is increased? 
        (b) How would you expect the conductivity of a semiconductor to change with temperature?

        \subsubsection{Answer}

    \subsection{Question 12}
        Three different materials have filled valence bands and empty conduction bands, and the Fermi energy lies in the middle of the gap. 
        The gap energies are 10 eV, 1 eV, and 0.01 eV. 
        Classify the electrical properties of these materials at room temperature and at 3 K.

        \subsubsection{Answer}

\section{Problem 7}
    Calculate the Coulomb energy and the repulsion energy for \ce{NaCl} at its equilibrium separation.

    \subsection{Solution}
        According to Table 11.1, \ce{NaCl} has an fcc structure.
        This means that the Madelung constant $\alpha = 1.7476$. 
        The coulombic potential energy of an ion in an ionic solid is defined by equation (\ref{eq:11.1}).
        \begin{equation}
            \label{eq:11.1} \tag{11.1}
            U_C = -\alpha \frac{e^2}{4\pi\varepsilon_0} \frac{1}{R}
        \end{equation}

        For \ce{NaCl}, the nearest-neighbor separation $R_0 = 0.281\,\unit{\nano\meter}$.
        Plug all of these into equation (\ref{eq:11.1}).
        \begin{align}
            U_C &=  -\alpha \frac{e^2}{4\pi\varepsilon_0} \frac{1}{R_0}
                =   -1.7476 \frac{(8.99\E{9}\,\unit{\frac{\newton\cdot\meter^2}{\coulomb^2}})(1.602\E{-19}\,\unit{\coulomb})^2}{0.281\E{-9}\,\unit{\meter}}\\
                &=  \boxed{-1.435\E{-18}\,\unit{\joule}}
        \end{align}

        You can arrive at the repulsion energy by dividing the Coulomb energy by the negative of the exponent $n$.
        In the case of \ce{NaCl}, $n = 9$ and $-n = -9$.
        \begin{equation}
            U_R =   \frac{U_C}{-n}
                =   \frac{-1.435\E{-18}\,\unit{\joule}}{-9}
                =   \boxed{1.594\E{-19}\,\unit{\joule}}
        \end{equation}

\section{Problem 11}

    \subsection{Solution}

\section{Problem 12}

    \subsection{Solution}

\section{Problem 13}

    \subsection{Solution}

\section{Problem 15}

    \subsection{Solution}

\section{Problem 16}

    \subsection{Solution}

\section{Problem 22}

    \subsection{Solution}

\section{Problem 23}

    \subsection{Solution}

\section{Problem 28}

    \subsection{Solution}

\section{Problem 29}

    \subsection{Solution}

\end{document}